% chapters/bagian1-pendahuluan-dan-latar-belakang-masalah.tex
% proposal.md bagian 1: Pendahuluan dan latar belakang masalah.
% Contoh teks ilmiah di bawah ini adalah placeholder untuk memperlihatkan
% hasil format; ganti dengan isi proposal yang sebenarnya.

\chapter{Pendahuluan dan Latar Belakang Masalah}

\section{Latar Belakang Masalah}
Perkembangan teknologi informasi telah mendorong lonjakan volume data
tekstual yang dihasilkan pengguna internet, terutama dalam bentuk
ulasan produk pada platform perdagangan elektronik. Ulasan tersebut
memuat opini konsumen yang sangat berharga bagi pelaku usaha, baik
untuk mengevaluasi kualitas produk maupun untuk merumuskan strategi
pemasaran. Namun demikian, jumlah ulasan yang terus bertambah setiap
hari membuat proses pembacaan dan penilaian opini secara manual
menjadi tidak efisien, sehingga dibutuhkan suatu pendekatan otomatis
yang mampu mengklasifikasikan sentimen ulasan tersebut ke dalam
kategori positif, negatif, atau netral.

Berbagai pendekatan pembelajaran mesin telah digunakan untuk
menyelesaikan permasalahan klasifikasi sentimen, mulai dari metode
statistik konvensional hingga metode pembelajaran mendalam (deep
learning) yang mampu menangkap hubungan kontekstual antar kata dalam
suatu kalimat \citep{example2022}. Salah satu arsitektur yang banyak
digunakan untuk data sekuensial seperti teks adalah Long Short-Term
Memory (LSTM), sebuah varian dari Recurrent Neural Network (RNN) yang
dirancang untuk mengatasi masalah \textit{vanishing gradient} pada
data berurutan yang panjang.

\subsection{Kesenjangan penelitian pada data berbahasa Indonesia}
Bahasa Indonesia memiliki karakteristik morfologis dan tata bahasa
yang berbeda dari bahasa Inggris, seperti penggunaan imbuhan, kata
tidak baku, dan campur kode dengan bahasa daerah maupun bahasa asing
pada ulasan produk. Karakteristik ini menjadi tantangan tersendiri
dalam praproses data sebelum data tersebut dapat digunakan untuk
melatih model klasifikasi \citep{example2023}. Oleh karena itu,
proposal penelitian ini mengusulkan penerapan metode LSTM untuk
klasifikasi sentimen ulasan produk berbahasa Indonesia, dengan
memperhatikan tahapan praproses yang sesuai dengan karakteristik
bahasa tersebut.
